\chapter{Áreas e volumes}\label{ch:g7:areas}

O Ano 6 mediu retângulos, \omterm{def:g6:shapes:triangles}{triângulos retângulos} e blocos
(\cref{ch:g6:measure}). Com uma tesoura --- corte um pedaço aqui, cole
ali --- as mesmas ideias dão agora as \omterm{def:g6:measure:area}{áreas} dos \omterm{def:g7:central:parallelogram}{paralelogramos}, de
\emph{todos} os triângulos e do disco, e os \omterm{def:g6:measure:volume}{volumes} dos \omterm{def:g7:areas:prism}{prismas} e dos
\omterm{def:g7:areas:prism}{cilindros}.

\section{Área de um paralelogramo}

\begin{theorem}[Paralelogramo]\label{thm:g7:areas:parallelogram}
Um \omterm{def:g7:central:parallelogram}{paralelogramo} com um lado de comprimento $b$ (a \emph{base} dele) e
distância $h$ entre esse lado e o oposto (a \emph{altura} dele) tem
\omterm{def:g6:measure:area}{área}
\[
A = b \times h .
\]
\end{theorem}

\begin{proof}[Demonstração com a tesoura]
Corte o \omterm{def:g6:shapes:triangles}{triângulo retângulo} que sobra de um lado e cole-o do outro: o
\omterm{def:g7:central:parallelogram}{paralelogramo} vira um retângulo $b \times h$ de mesma \omterm{def:g6:measure:area}{área}.
\end{proof}

\begin{omfigure}
\begin{tikzpicture}[scale=0.82]
  \fill[omDef!14] (0,0) -- (4,0) -- (5.2,1.8) -- (1.2,1.8) -- cycle;
  \draw[thick] (0,0) -- (4,0) -- (5.2,1.8) -- (1.2,1.8) -- cycle;
  \fill[omThm!25] (0,0) -- (1.2,1.8) -- (1.2,0) -- cycle;
  \draw[omThm, thick] (0,0) -- (1.2,1.8) -- (1.2,0) -- cycle;
  \draw[dashed] (1.2,0) -- (1.2,1.8);
  \node[below, font=\small] at (2.0,0) {$b$};
  \draw[Stealth-Stealth] (6.0,0) -- (6.0,1.8)
    node[midway, right, font=\small] {$h$};
  \begin{scope}[xshift=8.2cm]
  \fill[omDef!14] (0,0) rectangle (4,1.8);
  \draw[thick] (0,0) rectangle (4,1.8);
  \fill[omThm!25] (2.8,0) -- (4,1.8) -- (4,0) -- cycle;
  \draw[omThm, thick] (2.8,0) -- (4,1.8) -- (4,0) -- cycle;
  \node[below, font=\small] at (2.0,0) {$b$};
  \end{scope}
\end{tikzpicture}

{\small A demonstração com a tesoura: deslizar o triângulo vermelho da
ponta esquerda para a direita transforma o \omterm{def:g7:central:parallelogram}{paralelogramo} num retângulo
--- mesma \omterm{def:g6:measure:area}{área}, $b \times h$.}
\end{omfigure}

\begin{remark}
A altura é a distância \emph{perpendicular} entre as duas bases, e não
o comprimento do lado inclinado! Um \omterm{def:g7:central:parallelogram}{paralelogramo} bem inclinado pode
ter lados compridos e \omterm{def:g6:measure:area}{área} pequena.
\end{remark}

\section{Área de um triângulo}

\begin{theorem}[Triângulo]\label{thm:g7:areas:triangle}
Um triângulo de base $b$ e altura correspondente $h$ (a distância
perpendicular do \omterm{def:g5:solids:def}{vértice} oposto até a reta da base) tem \omterm{def:g6:measure:area}{área}
\[
A = \frac{b \times h}{2} .
\]
\end{theorem}

\begin{proof}
Duas cópias do triângulo, uma delas girada de meia volta em torno do
ponto médio de um lado, se encaixam num \omterm{def:g7:central:parallelogram}{paralelogramo} de base $b$ e
altura $h$ (é a \omterm{def:g7:central:def}{simetria central} em ação,
\cref{prop:g7:central:preserved}). O triângulo é a \omterm{def:g3:division:half}{metade} dele:
$\frac{bh}{2}$.
\end{proof}

\begin{omfigure}
\begin{tikzpicture}[scale=0.9]
  \fill[omDef!14] (0,0) -- (4,0) -- (1.4,2.1) -- cycle;
  \draw[thick] (0,0) -- (4,0) -- (1.4,2.1) -- cycle;
  \fill[omThm!12] (4,0) -- (1.4,2.1) -- (5.4,2.1) -- cycle;
  \draw[thick, omThm, dashed] (4,0) -- (5.4,2.1) -- (1.4,2.1);
  \draw[dashed] (1.4,0) -- (1.4,2.1);
  \draw[thin] (1.4,0) ++(-0.2,0) -- ++(0,0.2) -- ++(0.2,0);
  \node[below, font=\small] at (2,0) {$b$};
  \node[right, font=\small] at (1.42,1.0) {$h$};
\end{tikzpicture}

{\small Um triângulo e a cópia dele girada de meia volta (tracejada)
formam um \omterm{def:g7:central:parallelogram}{paralelogramo}: a \omterm{def:g6:measure:area}{área} do triângulo é
$\frac{b\times h}{2}$. Qualquer um dos três lados pode servir de base
--- com a \emph{altura correspondente a ele}.}
\end{omfigure}

\begin{example}\label{ex:g7:areas:triangle}
Um triângulo tem base de $9$ cm e a altura correspondente mede $4$ cm:
\[
A = \frac{9 \times 4}{2} = \frac{36}{2} = 18 \text{ cm}^2 .
\]
Num triângulo \emph{retângulo}, os dois catetos são uma base e a
altura dela: recuperamos a fórmula da
\cref{prop:g6:measure:formulas}.
\end{example}

\section{Área de um disco}

\begin{theorem}[Disco]\label{thm:g7:areas:disk}
Um disco de \omterm{def:g3:shapes:shapes}{raio} $r$ tem \omterm{def:g6:measure:area}{área}
\[
A = \pi r^2 .
\]
\end{theorem}

\begin{proof}[Ideia da demonstração]
Corte o disco em muitas fatias finas e iguais e disponha-as ponta com
ponta: elas quase formam um \omterm{def:g7:central:parallelogram}{paralelogramo} de altura $r$ cuja base é a
\omterm{def:g3:division:half}{metade} do comprimento da circunferência, $\pi r$
(\cref{prop:g6:measure:circle}). A \omterm{def:g6:measure:area}{área} dele fica perto de
$\pi r \times r = \pi r^2$, e a aproximação vira exata à medida que as
fatias ficam mais finas. Uma demonstração completa precisa do cálculo
integral, do volume do ensino médio.
\end{proof}

\begin{omfigure}
\begin{tikzpicture}[scale=0.85]
  \foreach \a in {0,45,...,315}
    \fill[omDef!20, draw=omDef] (0,0) -- (\a:1.4)
      arc (\a:\a+45:1.4) -- cycle;
  \begin{scope}[xshift=4.3cm, yshift=-0.15cm]
  \foreach \i in {0,1,2,3} {
    \fill[omDef!20, draw=omDef]
      ({\i*1.1},0) -- ({\i*1.1+0.55},1.4) -- ({\i*1.1+1.1},0) -- cycle;
    \fill[omDef!20, draw=omDef]
      ({\i*1.1+0.55},1.4) -- ({\i*1.1+1.1},0) -- ({\i*1.1+1.65},1.4)
      -- cycle;
  }
  \draw[Stealth-Stealth] (5.2,0) -- (5.2,1.4)
    node[midway, right, font=\small] {$\approx r$};
  \node[below, font=\small] at (2.5,-0.1) {base $\approx \pi r$};
  \end{scope}
\end{tikzpicture}

{\small Fatiar o disco e desenrolar as fatias: quase um \omterm{def:g7:central:parallelogram}{paralelogramo}
de base $\pi r$ (\omterm{def:g3:division:half}{metade} do contorno) e altura $r$.}
\end{omfigure}

\begin{example}\label{ex:g7:areas:disk}
O tampo redondo de uma mesa tem \omterm{def:g3:shapes:shapes}{raio} $60$ cm. A \omterm{def:g6:measure:area}{área} dele é
$\pi \times 60^2 = 3600\pi \approx 11\,310$ cm$^2$, um pouco mais que
$1.1$ m$^2$. Repare na \omterm{ex:g1:subtraction:difference}{diferença}: o \emph{\omterm{def:g6:measure:perimeter}{perímetro}} $2\pi r$ usa $r$,
e a \emph{\omterm{def:g6:measure:area}{área}} $\pi r^2$ usa $r^2$.
\end{example}

\section{Prismas e cilindros}

\begin{definition}[Prisma, cilindro]\label{def:g7:areas:prism}
Um \emph{prisma}\index{prisma} é um \omterm{def:g5:solids:def}{sólido} com duas \omterm{def:g5:solids:def}{faces} poligonais
\omterm{def:g6:lines:perp}{paralelas} e idênticas (as \emph{bases}), ligadas por retângulos; um
\emph{cilindro}\index{cilindro} é a mesma coisa com discos como bases.
A \emph{altura} é a distância entre as duas bases.
\end{definition}

\begin{theorem}[Volumes]\label{thm:g7:areas:volume}
Para um \omterm{def:g7:areas:prism}{prisma} ou um \omterm{def:g7:areas:prism}{cilindro} de \omterm{def:g6:measure:area}{área} da base $B$ e altura $h$:
\[
V = B \times h .
\]
\end{theorem}
\admitted

\begin{example}\label{ex:g7:areas:volume}
Uma barraca é um \omterm{def:g7:areas:prism}{prisma} deitado de lado: as bases dela são triângulos
de base $2$ m e altura $1.5$ m, e a barraca tem $3$ m de comprimento.
\begin{enumerate}
  \item \omterm{def:g6:measure:area}{Área} da base: $B = \frac{2 \times 1.5}{2} = 1.5$ m$^2$.
  \item \omterm{def:g6:measure:volume}{Volume}: $V = B \times h = 1.5 \times 3 = 4.5$ m$^3$.
\end{enumerate}
Uma lata de \omterm{def:g3:shapes:shapes}{raio} $4$ cm e altura $11$ cm:
$V = \pi \times 4^2 \times 11 = 176\pi \approx 553$ cm$^3$, mais ou
menos meio litro.
\end{example}

\section{Exercícios}

\begin{exercise}[$\star$]\label{exo:g7:areas:1}
Calcule a \omterm{def:g6:measure:area}{área} de um \omterm{def:g7:central:parallelogram}{paralelogramo} de base $8$ cm e altura $5$ cm; e a
de outro de base $6.5$ m e altura $4$ m.
\end{exercise}

\begin{exercise}[$\star$]\label{exo:g7:areas:2}
Um \omterm{def:g7:central:parallelogram}{paralelogramo} tem lados de $10$ cm e $6$ cm, e a altura relativa à
base de $10$ cm mede $4$ cm. Calcule a \omterm{def:g6:measure:area}{área} dele. Por que a resposta
não é $60$ cm$^2$?
\end{exercise}

\begin{exercise}[$\star$]\label{exo:g7:areas:3}
Calcule a \omterm{def:g6:measure:area}{área} de um triângulo de base $12$ cm e altura $7$ cm; de um
\omterm{def:g6:shapes:triangles}{triângulo retângulo} de catetos $9$ e $6$; e de um triângulo de base
$5.5$ m e altura $2$ m.
\end{exercise}

\begin{exercise}[$\star$]\label{exo:g7:areas:4}
Calcule a \omterm{def:g6:measure:area}{área} de um disco de \omterm{def:g3:shapes:shapes}{raio} $5$ cm e depois a de um semidisco de
\omterm{def:g3:shapes:shapes}{raio} $10$ cm (valores exatos com $\pi$ e depois arredondados à
unidade).
\end{exercise}

\begin{exercise}[$\star$]\label{exo:g7:areas:5}
Calcule o \omterm{def:g6:measure:perimeter}{perímetro} \emph{e} a \omterm{def:g6:measure:area}{área} de um disco de \omterm{def:g3:shapes:shapes}{raio} $3$ cm. Qual
das duas respostas vem em cm e qual em cm$^2$?
\end{exercise}

\begin{exercise}[$\star$]\label{exo:g7:areas:6}
Calcule o \omterm{def:g6:measure:volume}{volume} de um \omterm{def:g7:areas:prism}{prisma} de \omterm{def:g6:measure:area}{área} da base $24$ cm$^2$ e altura
$10$ cm; e o de um \omterm{def:g7:areas:prism}{cilindro} de \omterm{def:g3:shapes:shapes}{raio} $3$ cm e altura $8$ cm (exato e
depois arredondado).
\end{exercise}

\begin{exercise}[$\star$]\label{exo:g7:areas:7}
Um triângulo tem \omterm{def:g6:measure:area}{área} $28$ cm$^2$ e base $8$ cm. Encontre a altura
correspondente, escrevendo antes a equação.
\end{exercise}

\begin{exercise}[$\star\star$]\label{exo:g7:areas:8}
Trace um triângulo qualquer e meça com cuidado os três pares
base--altura. Calcule $\frac{b \times h}{2}$ para cada par: os três
resultados devem coincidir (a menos do erro de medição). Por quê?
\end{exercise}

\begin{exercise}[$\star\star$]\label{exo:g7:areas:9}
Um jardim em forma de trapézio pode ser dividido, por uma diagonal, em
dois triângulos que compartilham a mesma altura $h = 20$ m, com bases
$30$ m e $18$ m. Calcule a \omterm{def:g6:measure:area}{área} total dele. Você consegue adivinhar uma
fórmula geral para o trapézio?
\end{exercise}

\begin{exercise}[$\star\star$]\label{exo:g7:areas:10}
Um vaso cilíndrico de \omterm{def:g3:shapes:shapes}{raio} $6$ cm contém água até a altura de $15$ cm.
Toda a água é despejada numa caixa vazia de base
$18$ cm $\times$ $12$ cm. Que altura de água é atingida na caixa?
(Mesmo \omterm{def:g6:measure:volume}{volume}, base diferente.)
\end{exercise}

\begin{exercise}[$\star\star$]\label{exo:g7:areas:11}
Uma pizza de \omterm{def:g4:geometry:circle}{diâmetro} $30$ cm custa $9$ reais; uma de \omterm{def:g4:geometry:circle}{diâmetro} $40$ cm
custa $14$ reais. Calcule as duas \omterm{def:g6:measure:area}{áreas} e o preço por $100$ cm$^2$.
Que pizza vale mais a pena?
\end{exercise}

\begin{exercise}[$\star\star\star$]\label{exo:g7:areas:12}
Os dois catetos de um \omterm{def:g6:shapes:triangles}{triângulo retângulo} medem $6$ cm e $8$ cm, e a
hipotenusa mede $10$ cm. Calcule a \omterm{def:g6:measure:area}{área} dele com os catetos e depois
use essa \omterm{def:g6:measure:area}{área} para encontrar a altura relativa à \emph{hipotenusa}.
\end{exercise}

\section{Problema: o quadradinho que sumiu}

\begin{problem}[{Problema de fim de semana --- uma ``demonstração''
famosa de que $64 = 65$, desmascarada por áreas e inclinações, com
pizzas e anéis de sobremesa}]
\label{pb:g7:areas:1}
Um mágico corta um quadrado $8 \times 8$ em quatro peças, desliza-as e
as remonta num retângulo $5 \times 13$. A plateia conta:
$8 \times 8 = 64$, mas $5 \times 13 = 65$ --- uma unidade de \omterm{def:g6:measure:area}{área}
apareceu do nada! As fórmulas de \omterm{def:g6:measure:area}{área} deste capítulo são exatamente as
ferramentas que pegam o truque. Depois, o problema põe o raciocínio
honesto de \omterm{def:g6:measure:area}{áreas} para trabalhar: em trapézios, em pilhas inclinadas, na
economia da pizza --- e, por fim, volta ao mágico para um segundo
número, ainda mais estranho.

\medskip
\noindent\textbf{Parte I --- O truque, executado e desmascarado.}
As quatro peças do quadrado $8 \times 8$ são: dois \omterm{def:g6:shapes:triangles}{triângulos
retângulos} de catetos $8$ e $3$ e dois trapézios retângulos de lados
paralelos $5$ e $3$ e altura $5$.
\begin{enumerate}
  \item Calcule a \omterm{def:g6:measure:area}{área} do quadrado original e a \omterm{def:g6:measure:area}{área} do retângulo
        alegado. Enuncie o escândalo em uma linha.
  \item Calcule a \omterm{def:g6:measure:area}{área} de cada uma das quatro peças
        (\cref{thm:g7:areas:triangle}; para os trapézios, corte-os ou
        use o \cref{exo:g7:areas:9}).
  \item Some as quatro \omterm{def:g6:measure:area}{áreas}. Qual das duas figuras --- quadrado ou
        retângulo --- as peças conseguem de fato preencher?
  \item No arranjo do retângulo, a hipotenusa de um triângulo (que sobe
        $3$ ao longo de $8$) deveria continuar a borda inclinada de um
        trapézio (que sobe $2$ ao longo de $5$) numa única
        ``diagonal'' reta. Teste o alinhamento: $3$ ao longo de $8$ e
        $2$ ao longo de $5$ têm a mesma inclinação
        (\cref{thm:g7:prop:cross})?
  \item Explique onde se esconde o sexagésimo quinto quadradinho: que
        forma tem a falha ao longo da falsa diagonal, qual é a \omterm{def:g6:measure:area}{área}
        exata dela, e por que o olho não a percebe? (Qual é a largura
        média de uma lasca dessa \omterm{def:g6:measure:area}{área} esticada ao longo da diagonal do
        retângulo, de cerca de $13$ unidades?)
\end{enumerate}

\medskip
\noindent\textbf{Parte II --- Raciocínio honesto de \omterm{def:g6:measure:area}{áreas}.}
\begin{enumerate}[resume]
  \item Trace duas \omterm{def:g6:lines:perp}{retas paralelas} a $4$ cm de distância, uma base de
        $6$ cm sobre uma delas e três triângulos diferentes sobre essa
        base, com ápices em vários pontos da outra reta. Calcule as
        três \omterm{def:g6:measure:area}{áreas}. Por que elas são todas iguais, deslize o ápice
        onde deslizar sobre a reta dele?
  \item Demonstre a fórmula do trapézio em geral: um trapézio de lados
        paralelos $B$ e $b$ e altura $h$, cortado por uma diagonal, dá
        dois triângulos de mesma altura $h$. Conclua:
        \[
        A = \frac{(B + b) \times h}{2} .
        \]
  \item Aplique-a a um trapézio com $B = 9$ cm, $b = 5$ cm e
        $h = 4$ cm.
  \item Os marceneiros usam a mesma fórmula disfarçada: \omterm{def:g6:measure:area}{área} $=$
        (média dos dois lados paralelos) $\times$ altura. Explique por
        que essa é a mesma regra e refaça a questão 8 desse jeito.
  \item Incline uma pilha arrumada de cartas de baralho, deixando-a
        torta. Explique por que o \omterm{def:g6:measure:volume}{volume} não mudou e o que isso diz
        sobre a altura na fórmula do \omterm{def:g7:areas:prism}{prisma} $V = B \times h$
        (\cref{thm:g7:areas:volume}): que altura se deve tomar para uma
        pilha inclinada --- o comprimento inclinado ou o vertical?
        (Compare com o \cref{exo:g7:areas:2}, uma dimensão acima.)
\end{enumerate}

\medskip
\noindent\textbf{Parte III --- Anéis, pizzas e a volta do mágico.}
\begin{enumerate}[resume]
  \item Um lago circular de \omterm{def:g3:shapes:shapes}{raio} $3$ m fica no centro de um gramado
        circular de \omterm{def:g3:shapes:shapes}{raio} $5$ m. Calcule a \omterm{def:g6:measure:area}{área} do anel de grama (exata
        com $\pi$ e depois arredondada ao m$^2$).
  \item Encontre o \omterm{def:g3:shapes:shapes}{raio} do único disco cuja \omterm{def:g6:measure:area}{área} é exatamente igual à
        desse anel. (Os três \omterm{def:g3:shapes:shapes}{raios} que você tem agora formam um terno
        famoso, cuja história é contada no
        \cref{ch:g8:pythagoras}.)
  \item O que é mais pizza: uma pizza de \omterm{def:g4:geometry:circle}{diâmetro} $40$ cm ou duas
        pizzas de \omterm{def:g4:geometry:circle}{diâmetro} $28$ cm? Calcule e compare.
  \item Para \emph{dobrar} a \omterm{def:g6:measure:area}{área} de uma pizza, por que fator, mais ou
        menos, o \omterm{def:g3:shapes:shapes}{raio} dela tem de crescer? Teste o candidato $1.4$
        (lembre que as \omterm{def:g6:measure:area}{áreas} seguem o \emph{quadrado} do fator de
        escala, \cref{pb:g6:propdata:1}) e diga onde o fator exato ---
        o número cujo quadrado é $2$ --- faz a entrada oficial dele
        (\cref{ex:g8:pythagoras:diagonal}).
  \item O mágico volta: ele corta um quadrado $13 \times 13$ com os
        números $5$, $8$, $13$ fazendo os papéis que $3$, $5$, $8$
        faziam antes, e remonta um retângulo $8 \times 21$. Calcule as
        duas \omterm{def:g6:measure:area}{áreas}: o que some desta vez? Os números $3, 5, 8, 13, 21$
        são termos consecutivos da sequência do
        \cref{pb:g7:fractions:1} --- enuncie o padrão dos ganhos e
        perdas do mágico e a moral: por que as \omterm{def:g6:measure:area}{áreas} nunca mentem?
\end{enumerate}
\end{problem}
